% ============================================================================
% Tables for the RWR attacker sections -- EXTRACTED FROM THE LIVE DOC, 2026-08-22.
%
% NOTE: all three of these are ALREADY PRESENT in neurips_2026_v2_withtable.tex.
% This file is a clean copy for reference or re-paste. tab:primary's GCG rows and
% caption text are the other agent's work, restored after I briefly overwrote them. In
% particular tab:primary here INCLUDES the GCG rows added by the other agent --
% do not paste an older copy over it.
%
% Numbers come from:
%   probe_or/results/v5_judged/summary_v5.json     (prompted / reward-trained / control)
%   probe_or/results/judged_ablation_2model.json   (tab:ablation)
%   the GCG pipeline                               (gradient rows in tab:primary)
% ============================================================================


\begin{table}[t]
  \centering
  \small
  \caption{\textbf{Primary evaluation: how often does each method produce a genuine
  over-refusal?} All rows are scored by the same independent Claude Sonnet judge, which counts
  a rewrite only if it preserves the original request, remains benign, and the target refuses
  it; refusals that are justified by the rewritten text are discarded. The reward-trained row
  uses each model's best reward signal (the refusal-opener logit, Table~\ref{tab:reward_signal}).
  The prompted-attacker row is the \emph{same untrained model given the same instruction}, so it
  isolates the effect of training rather than of prompting. The control row is the unmodified
  benign prompts, establishing the floor an attacker must beat. Evaluated on 200 held-out
  originals, four rewrites each ($n=800$; $n=200$ for the control, which has no rewrites).
  The gradient row is not an attacker over held-out prompts but one optimization per prompt,
  so its denominator is every prompt the search ran on. It passes the same two-axis judge as
  the other rows, applied after the pipeline filter of Section~3.2. One asymmetry to note:
  the Llama filter includes a Llama Guard check on both prompts and the Qwen filter does not,
  so the Qwen rate is the less strictly gated of the two.}
  \label{tab:primary}
  \begin{tabular}{@{}llrrr@{}}
    \toprule
    Method & Target & Rewrites & Judge-confirmed & Rate \\
    & & evaluated & over-refusals & \\
    \midrule
    Gradient rewriting (GCG) & Llama-3-8B & $3{,}200$ & $1{,}057$ & $33.03\%$ \\
                             & Qwen3-32B  & $960$ & $256$ & $26.67\%$ \\
    \addlinespace[2pt]
    Prompted attacker & Llama-3-8B & $800$ & $39$ & $4.88\%$ \\
                      & Qwen3-32B  & $800$ & $23$ & $2.88\%$ \\
    \addlinespace[2pt]
    \textbf{Reward-trained attacker} & Llama-3-8B & $800$ & $\mathbf{80}$ & $\mathbf{10.00\%}$ \\
                                     & Qwen3-32B  & $800$ & $34$ & $4.25\%$ \\
    \addlinespace[2pt]
    Control (unmodified prompts) & Llama-3-8B & $200$ & $1$ & $0.50\%$ \\
                                 & Qwen3-32B  & $200$ & $0$ & $0.00\%$ \\
    \bottomrule
  \end{tabular}
\end{table}


\begin{table}[t]
  \centering
  \small
  \caption{\textbf{Which refusal signal makes the best RWR reward?}
  Every row is an attacker trained with the identical RWR recipe (LoRA $r{=}32$, 3 epochs,
  bin weights $[0,1,4,16]$, seed 42) on the same pool and the same held-out split; the
  \emph{only} difference between rows is the signal used to score candidate rewrites. The
  three vector layers were each selected by a different criterion: L17 by correlation with
  measured refusal change, L31 by classification AUC, and L12 by the abliteration test of
  \citet{arditi2024refusal}. Each trained attacker rewrites 200 held-out benign prompts, 4
  rewrites each ($n{=}800$); a two-axis Claude Sonnet judge counts a rewrite as over-refusal
  only if it preserves the original request, remains benign, and the target refuses it. The
  comparison column is the \emph{same base model, untrained}, given the same instruction on
  the same 200 prompts. Intervals are 95\% cluster bootstrap over the 200 originals
  (10{,}000 resamples), since the four rewrites of one prompt are not independent. Qwen has
  no abliteration-selected row because no single direction meaningfully reduces its
  harmful-prompt refusal, so there is no such layer to select.}
  \label{tab:reward_signal}
  \begin{tabular}{@{}llcccr@{}}
    \toprule
    Reward signal used to & Layer & \multicolumn{2}{c}{Judge-confirmed over-refusal} & & Diff. \\
    \cmidrule(lr){3-4}
    rank RWR training data & & trained attacker & untrained base & & (pp) \\
    \midrule
    \multicolumn{6}{@{}l}{\emph{Llama-3-8B-Instruct}} \\
    \quad Refusal vector    & 17 & $3.25$ {\tiny[1.75,\,5.00]} & $5.00$ {\tiny[3.38,\,6.75]} & & $-1.75$ \\
    \quad Refusal vector    & 31 & $3.62$ {\tiny[2.12,\,5.38]} & $3.12$ {\tiny[1.88,\,4.62]} & & $+0.50$ \\
    \quad Refusal vector    & 12 & $1.38$ {\tiny[0.50,\,2.50]} & $3.75$ {\tiny[2.25,\,5.38]} & & $-2.38$ \\
    \quad Multi-layer probe & all & $5.62$ {\tiny[3.50,\,8.12]} & $4.88$ {\tiny[3.12,\,6.88]} & & $+0.75$ \\
    \quad \textbf{Refusal-opener logit} & --- & $\mathbf{10.00}$ {\tiny[7.38,\,12.88]} & $4.88$ {\tiny[3.12,\,6.88]} & & $\mathbf{+5.12}^{*}$ \\
    \addlinespace[2pt]
    \multicolumn{6}{@{}l}{\emph{Qwen3-32B}} \\
    \quad Refusal vector & 58 & $2.00$ {\tiny[1.00,\,3.12]} & $1.50$ {\tiny[0.50,\,2.75]} & & $+0.50$ \\
    \quad Multi-layer probe & all & $4.50$ {\tiny[2.62,\,6.62]} & $2.50$ {\tiny[1.38,\,3.75]} & & $+2.00$ \\
    \quad Refusal-opener logit & --- & $4.25$ {\tiny[2.62,\,6.12]} & $2.88$ {\tiny[1.50,\,4.50]} & & $+1.38$ \\
    \bottomrule
  \end{tabular}
  \\[2pt]
  \footnotesize $^{*}$ the only row whose trained interval clears its own baseline interval.
\end{table}


\begin{table}[t]
  \centering
  \small
  \caption{\textbf{Removing a single direction reduces over-refusal without costing safety.}
  Each row ablates one direction at every layer and re-evaluates on two held-out sets: 400
  benign rewrites (over-refusal, which we want to fall) and 200 AdvBench harmful prompts
  (refusal, which must not fall). Both columns are percentage points of refusal \emph{removed}
  relative to the unablated model, so larger means more removed --- the goal on the left, the
  cost on the right. Refusal is scored by a Claude Sonnet judge on what the response does, not
  by matching a refusal phrase. \emph{No direction on either model moves harmful-prompt refusal
  by more than 1.0 points}, against random-direction nulls of 0.7 (Llama) and 1.5 (Qwen).
  Negative entries mean refusal went up. The final row gives the unablated rates the
  deltas are relative to, in their own columns.}
  \label{tab:ablation}
  \begin{tabular}{@{}lrrcrr@{}}
    \toprule
    & \multicolumn{2}{c}{Llama-3-8B-Instruct} & & \multicolumn{2}{c}{Qwen3-32B} \\
    \cmidrule(lr){2-3}\cmidrule(lr){5-6}
    Direction removed & over-refusal & harmful & & over-refusal & harmful \\
    & removed & refusal lost & & removed & refusal lost \\
    \midrule
    \textbf{overall direction} & $+6.6$ & $-1.0$ & & $+36.8$ & $+1.0$ \\
    whole-prompt refusal vector & $+28.7$ & $+0.6$ & & $+12.8$ & $-1.0$ \\
    \textbf{weaponization} & $+33.0$ & $-0.5$ & & $+3.9$ & $+0.0$ \\
    concealment & $+17.5$ & $-0.5$ & & $+15.7$ & $-0.5$ \\
    intrusion & $+18.1$ & $-0.5$ & & --- & --- \\
    exfiltration & $+7.8$ & $-0.5$ & & $+10.7$ & $-1.0$ \\
    exploitation & $+4.7$ & $+0.0$ & & $+9.4$ & $-0.5$ \\
    coercion & --- & --- & & $-1.8$ & $-0.5$ \\
    fabrication & --- & --- & & $-3.8$ & $-1.0$ \\
    \addlinespace[2pt]
    random direction (null) & $+0.7$ & $\sim0$ & & $+1.5$ & $\sim0$ \\
    \midrule
    \emph{unablated rate} & $74.2\%$ & $98.5\%$ & & $93.7\%$ & $99.0\%$ \\
    \bottomrule
  \end{tabular}
\end{table}
